\subsection*{同步实践：流水线、缓冲区与事件}

缓冲流水线有两项顺序要求：消费者必须等生产者填完槽位后才能读取，生产者也必须等消费者用完槽位后才能覆盖。装入再快也不能省去第二项要求，消费者再快也不能省去第一项要求。

第 3.6.2 节的双缓冲模式由一个组发起搬运、等待当前输入并轮换句柄。生产者—消费者流水线则将这些职责分配给不同异步区域，再用事件传递就绪和复用信息；在单组版本中，这些顺序由组内指令安排表达。

两种模式解决的是同一个根本问题：\textbf{不能在一方写入缓冲区的同时读取它}。区别只在于解法的组织方式。
